\section{子词分词：从有限词表到开放词表}

在课程的前几讲中，我们假设模型拥有一个\textbf{固定的词表} $V$。每个单词被映射到一个唯一的向量表示。然而，这种设计面临一个根本性问题：语言是\textbf{开放的}，新词、拼写变体、形态变化不断出现，固定词表无法覆盖所有可能的输入。

\begin{figure}[H]
\centering
\includegraphics[width=0.95\textwidth]{slides-images/slide-003.jpg}
\caption{固定词表的局限：常见词（如 pizza、tasty）有良好的嵌入，但拼写变体（如 taaaasty）、新造词（如 transformerify）只能被映射为 UNK 标记，丢失全部信息\protect\footnotemark}
\end{figure}
\footnotetext{Slides 第4页。}

\subsection{固定词表的问题}

当模型遇到训练时未见过的字符串时，只能将其映射为一个通用的 \texttt{<UNK>}（unknown）标记。这意味着：

\begin{itemize}
    \item 所有未知词共享同一个向量，\textbf{完全丢失了语义信息}
    \item 在形态丰富的语言（如斯瓦希里语、芬兰语、土耳其语）中，一个动词可能有数百种变形，词表膨胀极为严重
    \item 即使是英语，拼写变体、新造词、专有名词也层出不穷
\end{itemize}

\begin{figure}[H]
\centering
\includegraphics[width=0.95\textwidth]{slides-images/slide-005.jpg}
\caption{斯瓦希里语动词 \textit{ambia}（``to tell''）的变位表：超过 300 种变位形式。如果每种形式都占据词表一个位置，效率极低且难以学习共享的语义结构\protect\footnotemark}
\end{figure}
\footnotetext{Slides 第6页。来源：Wiktionary。}

\subsection{Byte Pair Encoding (BPE) 算法}

现代 NLP 的解决方案是\textbf{子词分词}（subword tokenization）。其核心思想是：不以``词''为最小单位，而是以\textbf{频繁共现的字符子串}为基本单元。

\begin{importantbox}{BPE 算法流程}
\begin{enumerate}
    \item 以所有单个字符（加上词尾标记）初始化词表
    \item 统计训练语料中所有相邻字符对的共现频率
    \item 将频率最高的字符对合并为一个新符号，加入词表
    \item 用新符号替换语料中的对应字符对
    \item 重复步骤 2--4，直到词表达到预设大小（通常 30k--50k）
\end{enumerate}
\end{importantbox}

\begin{figure}[H]
\centering
\includegraphics[width=0.95\textwidth]{slides-images/slide-007.jpg}
\caption{BPE 算法示意：从全字符词表开始，不断合并频繁共现的字符对，逐步构建子词词表\protect\footnotemark}
\end{figure}
\footnotetext{Slides 第8页。}

分词时，算法尽量用\textbf{最少的子词}来表示输入。例如：
\begin{itemize}
    \item 常见词 ``hat'' 直接作为一个 token
    \item 不常见的 ``taaaasty'' 被拆分为 ``ta'' + ``\#\#aa'' + ``\#\#sty'' 三个子词 token
    \item ``transformerify'' 可能被拆分为 ``transformer'' + ``\#\#ify''
\end{itemize}

其中 ``\#\#'' 前缀表示这个子词\textbf{不以空格开头}，即它是前一个子词的延续。

\begin{figure}[H]
\centering
\includegraphics[width=0.95\textwidth]{slides-images/slide-008.jpg}
\caption{子词分词的效果：常见词保持完整（如 hat、learn），未知词被拆分为已知子词序列，用三个有意义的嵌入替代一个无用的 UNK 嵌入\protect\footnotemark}
\end{figure}
\footnotetext{Slides 第9页。}

\begin{knowledgebox}{子词分词的关键性质}
\begin{itemize}
    \item 系统\textbf{不区分}一个 token 是完整的词还是子词片段——它们在模型中的地位完全相同，都是嵌入矩阵中的一个索引
    \item 如果需要获取一个被拆分的词的表示，可以\textbf{取平均}或取最后一个子词的上下文表示
    \item 现代 LLM（GPT、BERT、T5 等）几乎全部使用 BPE 或类似的子词分词方案
    \item 序列长度是 Transformer 的瓶颈，因此分词时倾向于用\textbf{更少更大}的子词
\end{itemize}
\end{knowledgebox}

\begin{warningbox}{子词分词不是``聪明''的语素切分}
BPE 是纯粹基于统计频率的算法，\textbf{不理解语言学意义上的词素}（morpheme）。它可能把 ``unhappiness'' 切分为 ``un'' + ``happiness''（恰好对应语素），也可能切分为 ``unh'' + ``appiness''（无语言学意义）。切分结果完全取决于训练语料的统计分布。
\end{warningbox}

\subsection{本章小结}

子词分词解决了固定词表的 UNK 问题，使模型能够处理任意输入字符串。BPE 通过迭代合并频繁共现字符对来构建词表，在词表大小和序列长度之间取得平衡。这一技术已成为所有现代语言模型的标配。

%% ========== Part 2: 从词嵌入到预训练 ==========

